\documentclass[preview]{standalone}

\usepackage{amsmath}
\usepackage{amssymb}
\usepackage{stellar}
\usepackage{definitions}

\begin{document}

\id{psicologia-coscienza}
\genpage

\section{Coscienza}

\begin{snippet}{coscienza-attenzione-introduzione}
    Coscienza e attenzione sono dimensioni fondamentali per l'uso di molte
    capacità psichiche, dalla percezione all'apprendimento. Pensieri non
    presenti nella consapevolezza possono giungere alla coscienza quando
    l'attenzione viene orientata verso di essi.
\end{snippet}

\begin{snippetdefinition}{coscienza-definition}{Coscienza}
    La \emph{coscienza} è la consapevolezza degli stimoli esterni e interni da
    parte del soggetto. È uno stato che accompagna un'esperienza risultante da
    processi di elaborazione delle informazioni.
\end{snippetdefinition}

\begin{snippet}{processi-inconsci-coscienza}
    Molte operazioni mentali si svolgono in modo inconscio. Solo in particolari
    condizioni, grazie all'intervento dei processi attentivi, le rappresentazioni
    prodotte dall'elaborazione delle informazioni diventano contenuti coscienti.
\end{snippet}

\begin{snippetdefinition}{conoscenze-dichiarative-definition}{Conoscenze dichiarative}
    Le \emph{conoscenze dichiarative}, o \emph{proposizionali}, sono conoscenze
    espresse sotto forma di proposizioni che stabiliscono relazioni tra idee e
    riguardano contenuti della vita quotidiana.
\end{snippetdefinition}

\begin{snippetdefinition}{conoscenze-procedurali-definition}{Conoscenze procedurali}
    Le \emph{conoscenze procedurali} sono conoscenze relative alle procedure con
    cui si svolgono compiti della vita quotidiana e dipendono dall'esercizio e
    dalla pratica.
\end{snippetdefinition}

\begin{snippetdefinition}{elaborazione-automatica-definition}{Elaborazione automatica}
    L'\emph{elaborazione automatica} è un'elaborazione rapida, poco dipendente
    dalle risorse attentive e generalmente non accompagnata da coscienza
    esplicita.
\end{snippetdefinition}

\begin{snippetdefinition}{elaborazione-controllata-definition}{Elaborazione controllata}
    L'\emph{elaborazione controllata} è un'elaborazione lenta, consapevole e
    dipendente dalle risorse attentive, che rende difficile svolgere altri
    compiti nello stesso momento.
\end{snippetdefinition}

\begin{snippet}{automatico-controllato}
    Non esiste un processo totalmente automatico: anche processi molto praticati
    possono richiedere di nuovo controllo cosciente in condizioni insolite. Ad
    esempio, guidare l'auto può avvenire in modo quasi automatico, ma torna a
    richiedere controllo se la strada è ghiacciata.
\end{snippet}

\begin{snippet}{james-coscienza}
    William James definisce la psicologia come descrizione e spiegazione degli
    stati di coscienza. Lo stato abituale di coscienza include percezioni,
    pensieri, sensazioni, immagini e desideri presenti in un dato momento, cioè
    l'attività mentale su cui è focalizzata l'attenzione.
\end{snippet}

\begin{snippetdefinition}{think-aloud-protocols-definition}{Think-aloud protocols}
    I \emph{think-aloud protocols}, o \emph{protocolli di verbalizzazione del
    pensiero}, sono resoconti prodotti dai partecipanti mentre svolgono un
    compito. Servono a documentare strategie mentali e rappresentazioni della
    conoscenza impiegate nello svolgimento del compito.
\end{snippetdefinition}

\begin{snippetdefinition}{experience-sampling-method-definition}{Experience-sampling method}
    L'\emph{experience-sampling method}, o \emph{ESM}, è un metodo di ricerca
    in cui i partecipanti forniscono informazioni sui propri pensieri e
    sentimenti in diversi momenti della giornata.
\end{snippetdefinition}

\begin{snippet}{funzioni-coscienza}
    La coscienza aiuta l'adattamento all'ambiente in tre modi:
    \begin{enumerate}
        \item svolge una \emph{funzione restrittiva}, riducendo il flusso di
        stimoli in ingresso e aiutando a ignorare le informazioni irrilevanti;
        \item permette l'\emph{immagazzinamento selettivo}, cioè la conservazione
        selettiva delle informazioni da analizzare;
        \item rende possibile la \emph{pianificazione}, permettendo di valutare
        alternative, usare conoscenze passate e immaginare conseguenze future.
    \end{enumerate}
\end{snippet}

\begin{snippetexample}{dilemma-morale-example}{Dilemma morale}
    Si immagini un aereo atterrato in emergenza in pieno oceano. C'è una sola
    scialuppa di salvataggio, ma sta affondando per il peso eccessivo. Una
    risposta immediata e guidata da processi inconsci può rifiutare il sacrificio
    di alcune persone; una valutazione consapevole e ponderata può invece
    concludere che il sacrificio salverebbe il maggior numero possibile di vite.
    L'esempio mostra che il comportamento può essere influenzato sia da processi
    inconsci sia da processi consci.
\end{snippetexample}

\end{document}
